\section{高斯信道：功率约束与连续容量}
\label{sec:07-Information-Theory/05-Channel-Coding/04}

实数加性白高斯噪声信道是最简单的连续信道模型：

\[
Y=X+Z,\qquad Z\sim\mathcal N(0,N),\qquad Z\perp X.
\]

你可能会想：只要把 \(X\) 的幅度放大，有效 SNR 就可以任意大——让不同码字在接收端隔得越来越远，不就能传任意多比特了吗？

这个推理只在忽略了一个物理事实时才成立：发送器的功率从来不是无限的。没有功率约束，容量公式没有有限上界——连续信道的容量必须与输入约束一起陈述才有意义。

\subsection{互信息上界导出容量公式}

最自然的约束是平均功率：

\[
\E[X^2]\le P.
\]

我们需要最大化输入分布下的互信息 \(I(X;Y)\)，其中 \(Y=X+Z\) 且噪声独立于输入。由互信息与微分熵的关系：

\[
I(X;Y)=h(Y)-h(Y\mid X)=h(Y)-h(Z).
\]

噪声 \(Z\sim\mathcal N(0,N)\) 的微分熵是固定的：

\[
h(Z)=\frac12\log_2(2\pi e N).
\]

因此最大化互信息等价于在输入功率约束下最大化输出微分熵 \(h(Y)\)。

输出 \(Y\) 的方差受功率约束间接限制。由 \(X\) 与 \(Z\) 独立：

\[
\Var(Y)=\Var(X)+\Var(Z)=\Var(X)+N.
\]

而方差不超过二阶矩：

\[
\Var(X)=\E[X^2]-(\E[X])^2\le\E[X^2]\le P.
\]

所以 \(\Var(Y)\le P+N\)，两个不等号在 \(\E[X]=0\) 且 \(\E[X^2]=P\) 时同时取等——即输入为零均值且功率用满。

现在关键的一步：给定方差上界的连续随机变量中，什么样的分布给出最大微分熵？答案是高斯分布。对于任意方差为 \(\sigma^2\) 的连续随机变量，

\[
h(Y)\le\frac12\log_2(2\pi e\sigma^2),
\]

等号当且仅当 \(Y\) 为高斯分布。所以

\[
h(Y)\le\frac12\log_2\bigl(2\pi e(P+N)\bigr).
\]

当输入取 \(X\sim\mathcal N(0,P)\) 时，\(Y=X+Z\) 是两个独立高斯的和，仍为高斯，方差恰好为 \(P+N\)，微分熵达到上界。代入互信息：

\[
I(X;Y)\le\frac12\log_2\bigl(2\pi e(P+N)\bigr)-\frac12\log_2(2\pi e N)
=\frac12\log_2\!\left(1+\frac{P}{N}\right).
\]

因此实标量 AWGN 信道的容量为

\[
C=\frac12\log_2(1+\mathrm{SNR})
\quad\text{bit/real use},
\]

其中 \(\mathrm{SNR}=P/N\) 是线性信噪比。高斯输入是达到容量的唯一连续分布。

这个推导将``高斯的微分熵最大''从一个看起来孤立的性质，变成了容量公式的直接原因：如果存在一个分布比高斯在相同方差下产生更大的微分熵，那就能超越上述容量——但这样的分布不存在。

\subsection{公式前面的 \(1/2\)、带宽与单位}

不同教材中 AWGN 容量公式的写法略有差异，原因在于计量口径。

\textbf{复基带信道。} 一个复符号等价于两个实维度（同相和正交分量）。若功率和噪声按复符号归一化，容量常写成

\[
C=\log_2(1+\mathrm{SNR})
\quad\text{bit/complex use},
\]

没有 \(1/2\)，因为已把两个实维度的容量合并。

\textbf{连续时间带限信道。} 带宽为 \(B\) Hz（双边）的信道每秒有约 \(2B\) 个实自由度。将每维容量乘以自由度数量：

\[
C=B\log_2\!\left(1+\frac{S}{N}\right)
\quad\text{bit/s},
\]

其中 \(S\) 是信号功率，\(N\) 是噪声功率（在带宽 \(B\) 内）。注意这里因子 \(1/2\) 被 \(2B\) 中的 2 消掉了——\(B\) 本身已经吸收了自由度数目。

每次使用、每个实维度、每个复符号和每秒，这四种计量口径不同。比较公式时不要只看有没有 \(1/2\)——先确认分母的单位是什么。

\subsection{数值直觉：dB 不是线性值}

SNR 通常用分贝（dB）报告，但公式里的 \(\mathrm{SNR}\) 必须是线性功率比，转换关系为

\[
\mathrm{SNR}_{\text{线性}}=10^{\mathrm{SNR}_{\text{dB}}/10}.
\]

把 dB 值直接代入 \(\log_2(1+\mathrm{SNR})\) 会得到一个荒谬的结果。

几个数值实例：

\begin{itemize}
\tightlist
\item
  \(\mathrm{SNR}=0\) dB \(\Rightarrow\) 线性值 \(=1\) \(\Rightarrow\) \(C=0.5\) bit/real use。即使信号和噪声一样强，每维仍能传半个比特。
\item
  \(\mathrm{SNR}=10\) dB \(\Rightarrow\) 线性值 \(=10\) \(\Rightarrow\) \(C\approx 1.73\) bit/real use。功率增加 10 倍，容量增加约 3.5 倍——远非线性。
\item
  \(\mathrm{SNR}=20\) dB \(\Rightarrow\) 线性值 \(=100\) \(\Rightarrow\) \(C\approx 3.33\) bit/real use。再增加 10 倍功率，容量只增加不到两倍。
\item
  \(\mathrm{SNR}=30\) dB \(\Rightarrow\) 线性值 \(=1000\) \(\Rightarrow\) \(C\approx 5.00\) bit/real use。
\end{itemize}

容量随 SNR 的增长在不同区间表现出截然不同的行为。

\textbf{低 SNR 区间（\(\mathrm{SNR}\ll1\)）：}

\[
\log_2(1+s)=\frac{\ln(1+s)}{\ln 2}\approx\frac{s}{\ln 2},
\]

容量近似线性依赖于 SNR。在噪声主导的区域，每增加一倍功率，容量几乎翻倍。噪声是瓶颈。

\textbf{高 SNR 区间（\(\mathrm{SNR}\gg1\)）：}

\[
\log_2(1+s)\approx\log_2 s,
\]

容量只以对数速率增长。在信号已经远强于噪声时，要多传 1 bit/real use，功率需要乘 4（约 6 dB）。功率效率急剧下降——这是香农极限的饱和效应。

这两种增长模式的转换意味着，在一个固定带宽的系统中投入更多功率，回报会递减。同时，在高 SNR 下用带宽换功率效率（扩频）是有效的——用更多自由度替代发送幅度。

\subsection{高斯星座与实际调制}

连续高斯输入达到容量，但实际发送器很少真的发送高斯分布的连续值。BPSK、QPSK、QAM 等有限星座的离散输入，其互信息在低 SNR 时接近容量曲线，但在高 SNR 时会饱和——星座只包含 \(M\) 个点，输入熵至多为 \(\log_2 M\)，无论 SNR 再高也不能超过这个上界。

这不表示容量公式``错了''。容量公式回答的是``在给定输入约束下，最优输入的极限是多少''。实际系统使用有限星座，意味着它解决的是一个带额外约束的问题：输入不仅受限于平均功率，还受限于离散星座结构、峰值幅度、功放线性区、DAC 量化精度。

如果问题变为``输入来自离散星座且功率受限''，最优输入分布就不再是连续高斯。如果再加上峰值幅度约束，达到容量的分布也不是高斯——对于幅度受限的实 AWGN 信道，容量达到分布是有限支撑的离散分布。每增加一层实际约束，容量公式都需要重新推导。

\subsection{容量是渐近界限，不是误码率公式}

\[
C=\frac12\log_2\!\left(1+\frac{P}{N}\right)
\]

的意义与离散信道容量完全相同：以低于 \(C\) 的速率，存在足够长的码使得错误概率任意小；高于 \(C\) 则无论怎么编码都不能可靠通信。

它不直接给出：
\begin{itemize}
\tightlist
\item
  未编码 QPSK 在某个 SNR 下的误比特率——那是调制方案的未编码性能，远低于容量界限；
\item
  给定块长和给定调制方式下可以达到的最大速率——那是有限块长信息论和编码调制的问题；
\item
  用有限精度实现的译码器的实际性能——容量不涉及硬件实现。
\end{itemize}

工程中做链路预算时，逐项核对以下问题能避免大部分数值错误：

\begin{itemize}
\tightlist
\item
  \(P\) 和 \(N\) 是否按相同的带宽归一化？
\item
  使用实维度还是复维度的单位？
\item
  SNR 是线性值还是已经转换？
\item
  约束包含的是平均功率、峰值功率，还是二者都有？
\item
  最终容量报告的是 bit/real use、bit/complex use 还是 bit/s？
\item
  接收端保留的是软信息（似然比）还是做了硬判决？有损硬判决会在容量链中引入不可逆的信息损失。
\end{itemize}

容量公式的简洁之下，每一处单位转换和约束定义都直接改变数值结果。
